\begin{table}[!t]
\centering
\caption{Standards and guidance documents compared by scope, treatment of business impact analysis, and nature of the requirements.}
\label{tab:lr_standards}
\footnotesize
\renewcommand{\arraystretch}{1.25}
\begin{tabularx}{\textwidth}{>{\raggedright\arraybackslash}p{2.6cm} >{\raggedright\arraybackslash}X >{\raggedright\arraybackslash}X >{\raggedright\arraybackslash}p{2.7cm}}
\toprule
\textbf{Document} & \textbf{Scope} & \textbf{Business impact analysis and recovery parameters} & \textbf{Nature of the text} \\
\midrule
ISO 22301:2019 \cite{iso22301catalogue} & Establishing and maintaining a business continuity management system suited to the organizational context. & Not stated in the published scope description. & Requirements standard. \\
ISO 22313:2020 \cite{iso22313catalogue} & Guidance on applying ISO 22301 to organizations of different types and contexts, so that delivery is sustained at a predefined acceptable capacity. & Not stated in the published scope description. & Guidance. \\
ISO/TS 22317:2021 \cite{iso22317catalogue} & Guidance on conducting a documented business impact analysis. & Adapts the analysis to organizational needs, objectives, resources, and constraints, and explicitly does not impose a single uniform process. & Technical specification (guidance). \\
NIST SP 800-34 Rev.~1 \cite{nist2010contingency} & Contingency planning for federal information systems; comprehensive business-process continuity is outside its scope. & Three analysis tasks (process criticality, resources, recovery priorities). Defines MTD, RTO, and RPO, and states that the RTO should normally be shorter than the MTD. Times and costs in the guide are illustrative. & Guidance with templates and examples. \\
\bottomrule
\end{tabularx}
\par\smallskip
\footnotesize Note: The entries for the ISO documents are based on their public catalogue descriptions.
\end{table}